\section{Unlocked access}
\label{sec:daring}

The dissertation opens with a simplification of Dekker's algorithm, which
a fence after each store makes correct on this machine:
\[
  \cseq{\cstore{f_1}{1}_1}{\cseq{\cfence_1}{\cload{r_1}{f_2}_1}}
  \;\oppar\;
  \cseq{\cstore{f_2}{1}_2}{\cseq{\cfence_2}{\cload{r_2}{f_1}_2}}
\]
From $f_1 = f_2 = 0$ the two registers should never both end at $0$. Each
thread reads a flag that the other writes, and neither takes the lock. The
calculus of Section~\ref{sec:calculus} cannot express this, because there
a thread may touch a cell only if it owns it or holds the lock.

\paragraph{The original rules.}
The dissertation's basic ``daring'' rules were meant for this case. They
let a thread load from or store to the invariant's cells without the lock,
provided that the single command, specified with $J$ in hand, gives $J$
back:
\[
  \inferrule{\spec{\femp}{J \fsep P}{a}{J \fsep Q}}{\spec{J}{P}{a}{Q}}
\]
The dissertation expected their soundness to be hard to show and left them
as future work. For the repaired calculus they are unsound. Adding the load rule lets
one derive a false specification
(\rocq{daring_load_extension_breaks_adequacy}), and each rule has an
instance whose premise is valid and whose conclusion is invalid
(\rocq{daring_load_honest_counterexample},
\rocq{daring_store_honest_counterexample}). Two features of the repaired
model are responsible. A thread that holds the lock also holds the
invariant's cells, so $J \fsep P$ can be false of a state the thread is
really in. And between lock operations the invariant's cells are fixed and
free of buffered writes, so a store without the lock can neither take the
cell it writes to nor leave its write there.

\paragraph{A restricted form.}
A second judgment, $\bspec{J}{P}{c}{Q}$, is defined for commands that
contain neither $\clock$ nor $\cunlock$. It differs from the first in
three ways.
\begin{itemize}[nosep]
  \item The invariant's resource may hold buffered writes, and $J$ may
    describe them with a history $e \fcell u \fseq e \fwrite{i} v$. A
    daring store places its write in the invariant. The dissertation noted
    that \rn{daring-store} requires this, and preferred invariants that
    describe the heap alone.
  \item Any primitive step may move cells between the invariant and the
    thread, where the first judgment allows this only at $\clock$ and
    $\cunlock$.
  \item Frames are disjoint from the thread's resource, and there is no
    temporal frame rule.
\end{itemize}
The rule for access is
\[
  \inferrule*[right=atomic]
    {\langle J \fsep P \rangle\; a\; \langle J \fsep Q \rangle}
    {\bspec{J}{P}{a}{Q}}
\]
for a single store, load, fence, assignment, $\cassume{}$ or
$\cassert{}$. Its premise is a statement about one machine step. From every
unlocked memory that splits into a model of $J$ and a model of $P$, the
step does not abort, and every result splits into a model of $J$ and a
model of $Q$. The invariant's part is read at full completeness, and the
thread's part gains no completeness mark except its own at a fence. No
condition is needed for the flushes that happen between steps, because
assertions are closed under flushing.

With \rn{atomic} the judgment has sixteen rules: the structural rules, the
spatial frame rule, the concurrency rule, and axioms for cells the thread
owns outright. It is sound and adequate for the machine
(\rocq{buffered_calculus_sound}, \rocq{buffered_derivation_adequacy}), with
no assumption on the number of processors.

\paragraph{The result.}
The program above is derivable with a postcondition that excludes
$r_1 = r_2 = 0$ (\rocq{buffered_dekker_derived};
Appendix~\ref{app:dekker} gives the invariant). Erasing the proof's token
cells, as for message passing, gives a theorem about the machine alone
(\rocq{dekker_fenced_all_runs}): from the heap $f_1 = f_2 = 0$ with empty
buffers, no execution aborts and no terminating execution ends with
$r_1 = 0$ and $r_2 = 0$. Without the fences such an execution exists
(\rocq{dekker_unfenced_zero_zero}).

\paragraph{Limits.}
This is the dissertation's simplified program, without the retry loop of
the full algorithm. The two
judgments are separate: a program cannot yet take the lock in one place and
use \rn{atomic} in another, and it is not known how to combine them. And a
thread needs a completeness mark before it can take a cell out of the
invariant, since a cell cannot be owned without one. A thread that starts
from $\femp$ may have none
(\rocq{buffered_emp_load_withdrawal_premise_invalid}). Starting each
thread with $\fbar{i}$, which costs nothing when the program begins,
supplies one (\rocq{buffered_barrier_parallel_initial_adequacy},
\rocq{buffered_marked_cell_withdrawal}).
